%% kapa: exam=1 round=27 year=2016 period=2 form=A subject=appraisal_law subject_ko=감정평가관계법규 seq=21 number=21
%% source: Q-Net 감정평가사 기출문제 (cst00309) artlSeq=5169402 file=extracted/1008/5169402/제27회 감정평가사 1차 문제/제27회 감정평가사 2교시문제(A형).pdf
%% conversion: kordoc 4.24.51 -> tools/gen1.py (auto); review=proofread-sonnet; flags=none
\begin{question}{21}
국유재산법령상 지식재산에 관한 설명으로 옳지 않은 것은?
\begin{choices}
  \item 「디자인보호법」에 따라 등록된 디자인권은 지식재산에 해당한다.
  \item 중앙관서의 장등은 지식재산의 사용허가등을 하려는 경우에는 수의(隨意)의 방법으로 할 수 있다.
  \item 상표권의 사용허가등의 기간은 10년 이내로 한다.
  \item 중앙관서의 장등은「중소기업기본법」에 따른 중소기업의 수출증진을 위하여 필요하다고 인정하는 경우 지식재산의 사용허가에 따른 사용료를 면제할 수 있다.
  \item 저작권등의 사용허가등을 받은 자는 해당 지식재산을 관리하는 중앙관서의 장등의 승인을 받아 그 저작물을 변형할 수 있다.
\end{choices}
\end{question}
\answer{3}
